\chapter{樹状突起の数によるニューロンの分類}
\chaptermark{樹状突起の数}
\label{sec:dendrites}

\section{回路パラメータとしての入力数}

第~\ref{sec:neurontypes}~章ではニューロンを出力が放出するものによって、第~\ref{sec:eetypes}~章ではどんな素子として働くかによって分類した。エンジニアがすぐに気づく第三の特徴もある。\emph{ニューロンの入力はいくつあるか}である。入力を集めるのは樹状突起、すなわち他のニューロンの軸索が接続する枝である（第~\ref{sec:dofa}~章、第~\ref{sec:weightstore}~節）。樹状突起がまったくないニューロンもあれば、1本か2本のものも、10万個の入力を集める木をもつものもある。電子技術者の言葉ではこれはファンインであり、ほとんどの場合、課題に見合っている。

樹状突起の数と大きさがなぜそれほど重要かは、簡単な見積もりからわかる。膜は比容量$c_m\approx1$~\textmu F/cm$^2$のコンデンサであり、ニューロンと樹状突起を合わせた面積$A$が大きいほど、電圧を$\Delta V$上げるのに運ぶべき電荷が多い。漏れを無視すれば、入力電流$I$がこれを行う時間は
\begin{empheq}[box=\keybox]{equation}
t \;\approx\; \frac{c_m\,A\,\Delta V}{I} .
\label{eq:dendriteload}
\end{empheq}
である。短い樹状突起が数本の小さなニューロンは、面積が約$2\cdot10^{-5}$~cm$^2$、容量が約$20$~pFである。数nAの電流を流す1個の大きなシナプスは、これを0.1ミリ秒未満で$20\mV$持ち上げる。大きな木をもつニューロンの面積は約15倍で、容量は約$300$~pFである。約$0.1$~nAの普通のシナプスなら充電に$60\ms$かかる。これは漏れの時定数よりはるかに長いので、実際には決して充電できない。このようなニューロンには数十個の同時入力が必要である。ここでの数値は桁の見積もりだが、結論はそれに左右されない。\emph{入力が少なければ速く正確、多ければ遅く総和で働く}。

\begin{figure}[t]
\centering
\begin{tikzpicture}[>=Stealth,
  soma/.style={circle,draw,very thick,fill=blue!6,minimum size=0.55cm,inner sep=0pt},
  ax/.style={->,very thick,red!70!black},
  den/.style={very thick,blue!60!black}]
% 0 dendrites: sensor on a line
\begin{scope}[xshift=0cm]
\draw[den,decorate,decoration={snake,amplitude=1pt,segment length=4pt}] (-0.9,0) -- (0,0);
\node[font=\scriptsize] at (-0.9,0.3) {センサ};
\draw[ax] (0,0) -- (1.0,0);
\draw[thick] (0.1,0) -- (0.1,0.45);
\node[soma,minimum size=0.4cm] at (0.1,0.65) {};
\node[font=\scriptsize,align=center] at (0.05,-0.75) {$0$\\センサ付き線路};
\end{scope}
% 1 dendrite
\begin{scope}[xshift=3.3cm]
\draw[den] (-0.9,0) -- (-0.28,0);
\node[soma] at (0,0) {};
\draw[ax] (0.28,0) -- (0.9,0);
\node[font=\scriptsize,align=center] at (0,-0.75) {$1$\\バッファ、中継器};
\end{scope}
% 2 dendrites: one per ear
\begin{scope}[xshift=6.2cm]
\draw[den] (-0.28,0) -- (-0.9,0); \draw[den] (0.28,0) -- (0.9,0);
\node[soma] at (0,0) {};
\draw[ax] (0,-0.28) -- (0,-0.6);
\node[font=\scriptsize] at (-0.9,0.25) {左}; \node[font=\scriptsize] at (0.9,0.25) {右};
\node[font=\scriptsize,align=center] at (0,-1.0) {$2$\\時間のAND};
\end{scope}
% ~4 short dendrites: mixer
\begin{scope}[xshift=9.0cm]
\foreach \a in {45,135,225,315}{\draw[den] (\a:0.28) -- (\a:0.7); \fill[blue!60!black] (\a:0.72) circle (0.06);}
\node[soma] at (0,0) {};
\draw[ax] (0.28,0) -- (1.0,0);
\node[font=\scriptsize,align=center] at (0,-1.0) {$\sim4$\\ミキサ};
\end{scope}
% many: summing tree
\begin{scope}[xshift=12.0cm]
\foreach \a in {60,90,120}{\draw[den] (0,0.28) -- ++(\a:0.55) coordinate (b); \foreach \d in {-20,20}{\draw[den] (b) -- ++(\a+\d:0.35);}}
\foreach \a in {-120,-60}{\draw[den] (0,-0.28) -- ++(\a:0.45);}
\node[soma] at (0,0) {};
\draw[ax] (0.28,0) -- (1.0,0);
\node[font=\scriptsize,align=center] at (0,-1.0) {$10^4$--$10^5$\\加算器};
\end{scope}
\end{tikzpicture}
\caption{入力数による五つのクラス。青は樹状突起（入力）、赤は軸索（出力）。入力が少ないほど、ニューロンは個々の信号を速く正確に伝える。多いほど、足し合わせと分類が得意になる。細胞の形の図ではなく模式図である。}
\label{fig:dendriteclasses}
\end{figure}

\section{樹状突起ゼロ：線路の端のセンサ}

触覚、痛み、筋伸張の感覚ニューロン（第~\ref{sec:sensors}、\ref{sec:pain}、\ref{sec:muscles}~章）には、細胞体に樹状突起が1本もない。軸索は細胞体から1本の小枝として出て、すぐに二つに分かれる。一方の枝は皮膚や筋肉へ向かい、その末端自体がセンサになる。もう一方は脊髄へ向かう。スパイクはセンサから細胞体を経ずにまっすぐ脊髄へ走る。エンジニアの目には、遠端にセンサがある通信線路であり、細胞体は電源ユニットのように脇にぶら下がっている。線路に給電するが、伝送には加わらない。利点は速さと信頼性である。経路上に、信号を足し合わせて送るかどうかを改めて判断する場所が一つもない。

光と音のセンサはさらに極端な例である。これらは入力を集めるニューロンではなく変換器そのものであり、入力はシナプスではなく受容体タンパク質である（図~\ref{fig:transducer}）。

\section{樹状突起1本：バッファ}

樹状突起1本と軸索1本をもつニューロンは、1本の入力線を受け取って先へ渡す。嗅覚ニューロンがそうである。唯一の樹状突起が鼻の表面へ伸び、ただ1種類の受容体をもつ~\cite{Mombaerts2004}。光センサと眼の出力ニューロンの間にある網膜の中継細胞も、耳のセンサと脳をつなぐニューロンもそうである。エンジニアならバッファ、あるいは中継器と呼ぶだろう。計算はせず、1個の信号を届ける。ときには形を変える。たとえば第~\ref{sec:sensors}~章のように、センサの滑らかな電圧をスパイクに変える。

\section{樹状突起2本：時間のAND}

音の方向を決める同時検出器（図~\ref{fig:itd}）は、哺乳類ではしばしば反対方向を向いたちょうど2本の樹状突起をもつ。一方には左耳からの入力が、他方には右耳からの入力が届く~\cite{Grothe2010}。なぜ入力を分けるのか。式~\eqref{eq:shunt}を思い出そう。膜の一か所で多くの興奮性チャネルが開いていると、そこの電圧は平衡電位に近づき、入力が一つ増えるごとに加わる分は小さくなる。したがって、片耳からの多くの入力を1本の樹状突起に集めるとそれは飽和し、それだけではニューロンを閾値に届かせられない。一方、各樹状突起からの1個ずつの入力はほぼ線形に足し合わされる。2本の樹状突起はニューロンをより厳密なANDに変える。信号が両側から同時に来たときだけ発火せよ、である。モデルでは、このように入力を分けることが実際に同時検出を改善することが示されている~\cite{AgmonSnir1998}。

\section{短い樹状突起が数本：ミキサと中継器}

\emph{ミキサ。}第~\ref{sec:motorlearning}~章の運動学習の回路では、約700億個の小さな\nt{GLUT}ニューロンが入力を非常に長いベクトルに拡張する。それぞれの樹状突起は短いものが約4本だけで、各々の先端は1個の入力の終末をつかむ「爪」になっている。つまり各ミキサは多くの入力のうち約4個しか見ず、その4個に対する小さなAND素子として働く。$M$本の入力線から選べる4個組は$\binom{M}{4}$通りあり、$M=100$ならほぼ400万である。なぜそれほど必要なのか。$n$入力の閾値素子が二つのクラスに正しく分けられるランダムなパターンは、たかだか
\begin{empheq}[box=\keybox]{equation}
P_{\max} \;\approx\; 2n
\label{eq:cover}
\end{empheq}
個である~\cite{Cover1965}。同じ回路の出力ニューロンはミキサから約$10^5$個の入力を受けるので、\eqref{eq:cover}によれば上限は約$2\cdot10^5$個の異なる対応である。これは理想的な素子の限界である。興奮性シナプスのように重みがすべて正で、ノイズへの余裕も必要なら、同じニューロンの見積もりは約$4\cdot10^4$個の対応になる~\cite{Brunel2004}。入力の少ないミキサが必要なのは、まさに大きな加算器に多数の異なる、ほぼ独立な入力を与えるためである~\cite{Marr1969,Albus1971}。

\emph{中継器。}耳から方向検出器への経路には、短い樹状突起を少数もつニューロンがある。これらは巨大なシナプスをわずか数個、あるものは実質的に1個だけ受け取る。\eqref{eq:dendriteload}によれば、これはまさに必要なものである。小さな容量と大きな電流により、ニューロンは入力スパイクごとに1ミリ秒未満で発火し、時間精度をほとんど失わない。そうした中継器の一つは\nt{GLUT}を受け取り\nt{GLYC}を出す。数十マイクロ秒の精度をもつインバータであり、方向検出器に速い抑制を供給する~\cite{BorstSoria2012}。

\section{数千の入力：加算器と分類器}

尺度の反対の端には、1万個の入力をもつ皮質の\nt{GLUT}ニューロンと、10万個の入力をもつ運動学習回路の出力\nt{GABA}ニューロンがある。\eqref{eq:dendriteload}によれば、こうしたニューロンは遅く、個々の入力には応答できない。足し合わせるのである。これは第~\ref{sec:glutcomputer}~章の積分ニューロンであり、分類器を組み立てる加算器である。大きな木は新しいものも加える。その枝はそれぞれ閾値をもつ別々のサブ回路として働くので、1個のニューロンが小さな多層ネットワークのように振る舞う~\cite{Beniaguev2021,Gidon2020}。

\section{出力でもある樹状突起}

軸索がまったくなく、樹状突起が入力と出力の両方を兼ねるニューロンがある。信号を受け取り、自ら神経伝達物質を放出する。最もよく調べられた例は、運動の方向の検出にかかわる網膜の\nt{GABA}ニューロンである（第~\ref{sec:gaba}~章、図~\ref{fig:direction}）。このニューロンでは、各枝が単独で細胞体から自分の先端へ向かう動きに応答し、その先端から抑制を出す~\cite{Euler2002}。1個のニューロンが、枝ごとに一つずつ、複数の独立した方向検出器になっている。エンジニアから見れば、これは1個の素子ではなく、1個のパッケージに複数のチャネルを収めたICである。

\section{まとめ}

\begin{table}[t]
\centering
\footnotesize
\setlength{\tabcolsep}{3pt}
\renewcommand{\arraystretch}{1.15}
\begin{tabular}{>{\raggedright}p{1.9cm}>{\raggedright}p{3.4cm}>{\raggedright}p{2.6cm}>{\raggedright}p{1.8cm}>{\raggedright\arraybackslash}p{3.8cm}}
\toprule
樹状突起 & 例 & 素子 & 入力数 & わかっていること \\
\midrule
$0$ & 触覚、痛み、伸張のセンサ & 端にセンサのある線路 & --- & 確立 \\
$1$ & 嗅覚ニューロン、網膜の中継細胞、耳のニューロン & バッファ、中継器 & $1$--少数 & 確立 \\
$2$ & 音の方向検出器 & 時間のAND & 2束 & 構造は確立。入力を分ける利点はモデルで示された \\
$\sim4$ & 運動学習回路のミキサ & 小さなAND素子 & $\sim4$ & 確立。入力拡張の役割は根拠の強い仮説 \\
少数、短い & 耳からの経路の中継器 & 中継器、インバータ & $1$--数個の巨大シナプス & 確立 \\
数千 & 皮質の\nt{GLUT}ニューロン、運動学習の出力ニューロン & 加算器、分類器 & $10^4$--$10^5$ & 確立。サブ回路としての枝は個別の例で測定 \\
出力を兼ねる樹状突起 & 網膜の方向\nt{GABA}ニューロン & 多チャネルIC & 多数 & 測定済み \\
\bottomrule
\end{tabular}
\caption{樹状突起の数によるニューロンの分類。}
\label{tab:dendrites}
\end{table}

\begin{keyblock}
\begin{proposition}[入力数は課題に従う]
\label{prop:dendrites}
個々の信号の速さと精度が必要な場所、すなわちセンサ、中継器、同時検出器では、ニューロンの樹状突起は少なく、入力も少なく、シナプスは大きい。小さな容量~\eqref{eq:dendriteload}はすばやく充電される。足し合わせと分類が必要な場所では、ニューロンは数千の入力をもつ大きな木をもつ。クラスの構造は測定済みである。計算上の説明は根拠が強いが、多くのクラスについては仮説のままである。
\end{proposition}
\end{keyblock}

表~\ref{tab:dendrites}は、これまでの二つの分類、すなわち神経伝達物質による分類（表~\ref{tab:types}）と素子による分類（表~\ref{tab:eetypes}）を補う。三つとも同じ素子を別々の側から見ている。何を出すか、何を計算するか、どれだけ聞くかである。

\section{問題}

\begin{exercise}[\lvlA]
\label{ex:19:1}
\emph{大きなニューロンに入力はいくつ必要か。}$C=300$~pF、$\tau_m=20\ms$のニューロンが、それぞれ$0.1$~nAの電流源であるシナプスを受け、閾値は静止電位より$20\mV$高い。(a)~そもそも閾値に達するには、$10\ms$で達するには、$5\ms$で達するには、同時に何個のシナプスが働く必要があるか。(b)~$C=20$~pFのニューロンは$4$~nAのシナプス1個でどれだけの時間で閾値に達するか。
\end{exercise}

\begin{exercise}[\lvlA]
\label{ex:19:2}
\emph{なぜ4なのか。}ミキサは$M=100$本から選んだ$k$本の線のANDで、パターンでは半数の線が活動し、ミキサは$7\cdot10^{10}$個ある。(a)~$k=2$、$4$、$40$でパターンに応答するミキサは何個か。(b)~$10^5$入力の出力ニューロンについて、ミキサは$M=100$本の直接の線に比べ対応の数~\eqref{eq:cover}を何倍にするか。
\end{exercise}

\begin{exercise}[\lvlB]
\label{ex:19:3}
\emph{$2n$はどこから来るか。}原点を通る超平面は、$n$次元空間の一般の位置にある$P$個の点を$C(P,n)=2\sum_{k=0}^{n-1}\binom{P-1}{k}$通りに二つのクラスに分ける。(a)~$P=2n$で$2^P$通りの分け方のちょうど半分が分離可能であることを示せ。(b)~$n=100$、$P=180$と$220$で分離可能な割合はいくらか。
\end{exercise}

\begin{exercise}[\lvlB]
\label{ex:19:4}
\emph{厳密なANDとしての2本の樹状突起。}樹状突起は漏れ$g_L$をもつ区画で、各入力はそこに電位$E_E$のコンダクタンス$g_0=0.5\,g_L$を開く。細胞体は$\kappa(V_1+V_2)$を見て、閾値は$\theta=1.1\,\kappa E_E$である。なぜ1本の樹状突起にいくつ入力があってもニューロンは発火しないのか。各樹状突起にいくつの入力が必要か。
\end{exercise}

\begin{exercise}[\lvlB]
\label{ex:19:5}
\emph{中継器の設計。}スパイクのジッタは$\sigma_t=\sigma_V/\dot V$で、$\dot V$は閾値付近での電圧の傾きである。ノイズ$\sigma_V=1\mV$で$\sigma_t\le20$~\textmu sが必要とし、漏れは無視する。$C=300$~pFのニューロンには$0.1$~nAの同期したシナプスが何個必要か。$C=20$~pFで$4$~nAのシナプス1個のニューロンのジッタはいくらか。
\end{exercise}

\begin{exercise}[\lvlB]
\label{ex:19:6}
\emph{シミュレーション：長い樹状突起の充電。}長さ$L=\ell/\lambda$、両端が閉じた受動ケーブル$\tau_m\partial_tV=\lambda^2\partial_x^2V-V+i/g$をシミュレートせよ（$40$区画、オイラー法）。遠端に短い電流パルスを入れる。$L=0.2$、$1$、$2$のとき、近端はいつ、同じ電荷を樹状突起全体に広げたときの電圧の何割まで上がるか。
\end{exercise}

\begin{exercise}[\lvlC]
\label{ex:19:7}
\emph{研究：最適な入力数。}容量$C=C_0+nc_s$、同時に働く入力は電流$I_s$のものが$m$個、ニューロンは$P\approx2n$個の対応を記憶する~\eqref{eq:cover}。$m$一定の場合と割合一定$m=\varphi n$の場合に、応答時間は$P$にどう依存するか。コスト基準を提案し、中継器と分類器それぞれの最適な$n$を求めよ。
\end{exercise}
